% ============================================================
% Part III — Desktop Section — M5 (Design — The Levers as Archai) —
%   DRAFT v1
% Mode: Stage 3 guided drafting (per-paragraph protocol). Plan approved
%   2026-07-28. Band G; style canon = approved M3.1 sample; conventions
%   as DESKTOP-M3/M4 drafts (MLA; respondent-code parentheticals; chain
%   nodes A\textsubscript{n}; no module labels or hypothesis notation in
%   rendered text; no metatextual forward-refs; no stylistic
%   metalanguage; NEVER sentence-initial "And"; NEVER "record" as name
%   for evidence; % symbol per-student percentages; student typos
%   verbatim, no [sic]; "YouTube platform/route"; terminology lock;
%   "being-in" hyphenated).
%
% VOCABULARY RULING (user 2026-07-28): keep technical terms that do
%   analytical work — coordination game, equilibrium, activation
%   energy — REMOVE all cost/price/budget/spend metaphor language per
%   the standing ban; chronos = a DESIGN MATERIAL (allotted/reserved,
%   never spent/budgeted). User: "we can and will revise as we go."
%
% Architecture (approved): M5.1 (2 ¶¶ levers-as-archai) → M5.2 (5 ¶¶,
%   ★ the expanded WE lever: equilibrium · synchrony constraint ·
%   rung 0 + low ladder · ladder top w/ ignition logic + bia guard ·
%   motto) → M5.3 (1 ¶ motto/mirror w/ capture-economy breath).
% INHERITANCE: M4.3 (hold-open-vs-convene; floors); M3.5 (locked spec);
%   M0.2 (systemwide asynchrony); M1.6 locked bank (Kreijns Pitfall 2
%   full deployment HERE; De Back inverse group-size; Garrison binding-
%   element; van der Meer verbal-audio-essential); M3.3 ¶5 (bia/hexis
%   guard). Anchors: cross-case §3.3; deployment-brief §5; F7 correction
%   + recount.
% ============================================================
% --- M5.1 The levers, ranked, as archai (¶¶1–2 DRAFTED 2026-07-28 — in
%     review. ARCHĒ LOCUS PINNED: Part I §1.1-v3 footnote — six senses at
%     Metaphysics V.1 (Δ.1) 1012b34–1013a23, chain uses the FOURTH
%     ("that from which...a thing first arises, and from which the
%     movement or the change naturally first proceeds," 1013a7–10, Ross
%     wording as quoted in Part I; outline's "Part I §3" pointer was
%     off — it is §1.1). Lever content per deployment-brief §5 +
%     cross-case §3.3; threshold "quarter of the class" inherits M3.1;
%     orekton lever inherits S25-R035 thread (no re-quote); pacing lever
%     kept AFTERMATH-FREE (depletion/carryover = lab-pole evidence,
%     reference only — phrased as design-outcome claim instead).
%     EXPANSION CANDIDATE held: ASEE published planned fixes
%     (MacOS/Linux, reduced storage, mobile — vle-05 §5) could cite into
%     the threshold lever; page unpinned, add on user order.
%     P2 RULINGS APPLIED 2026-07-28 — ¶2 REWRITTEN as ¶¶2–3:
%     (1) plain-named lever list up front: THE DOOR (grounded in M3.1's
%     own "only door left open" image) · THE TASK (chain A4) · THE FIND
%     (needs-finding + orekton) · THE MEETING (M4.3's own term) · THE
%     WHILE (FCM Weile — "the same while that boredom stretches");
%     (2) "threshold" dropped for door language, made concrete;
%     (3) task lever slowed: outside-completable assignments → in-world
%     work → reason-to-enter → justification answer;
%     (4) find lever recast as OBJECT OF DESIRE (not of the look):
%     mandate-context curation — name what a needs-finder watches for,
%     caption/transcribe, deliverable-linked finds, finds that reward a
%     second pair of eyes (desirable object doubles as occasion to meet
%     → feeds the meeting lever);
%     (5) "co-pole" removed — POLE BAN surfaced and recorded
%     ([[feedback-no-pole-language]]); meeting lever grounded in
%     co-disclosedness. RETROACTIVE SWEEP of M4.3 executed same day
%     (6 instances: register/two sides/world-disclosedness maximized/
%     "Co-disclosedness it did not carry"/neither in kind/rests on the
%     thinnest supply). OPEN FLAG: M3.6 approved prose has ONE "both
%     poles" in the divergence-pair sense ("One voice, both poles") —
%     user ruling needed before touching the M3 file.
%     P2/P3 v2 RULINGS APPLIED 2026-07-28 — ¶¶2–3 REWRITTEN as ¶¶2–4
%     (M5.1 now 4 ¶¶): (P2.1) lever names ITALICIZED throughout;
%     (P2.2) third-term question FOOTNOTED w/ W26-R033/Q9 + W26-R056/Q9
%     (ledger reuses); (P2.3) TASK-LEVER DOCTRINE CORRECTED (AUTHOR):
%     in-world requirement = compulsion = "pedagogically prosaic," never
%     the justification; the world earns existence by being BETTER —
%     rhetorical incorporation as the mechanism of pedagogical
% TASK-LEVER QUOTATION PASS 2026-07-28 (user-approved set; Makransky
%   genetics HELD OUT by user ruling). Backup 20260728T155206.
%   Insertion 1 (¶3): Dalgarno & Lee 2010 [VoR year] — affordance
%     sentence + trivial-fact warning, BOTH PDF-verified char-exact,
%     TRUE p. 25 (entry loci ran 1 high; form-feed count = 23 pp =
%     10–32 exact); W26-R033/Q9 fragment re-verified vs CSV line 787.
%   Insertion 2 (new ¶ after ¶4 close): Cook et al. 2011 JAMA meta —
%     609 studies / 35,226 trainees / "question the need" quote
%     PDF-verified TRUE p. 987; scoped honestly (vs.-no-intervention,
%     medium-agnostic, "without ranking routes"). Dubovi, Levy & Dagan
%     2017 — PILL-VR desktop clinical sim; gains vs lecture held at
%     5-month retest (23); 9/10-vs-9/10 dissociation SCOPED to n=10
%     recorded subsample (23–24); control + presence-not-unique quotes
%     PDF-verified char-exact TRUE p. 25 (entry said 26 — 1 high;
%     form-feed count = 12 pp = 16–27 exact).
%   Works Cited adds queued: Dalgarno & Lee 2010 (short title iff 2012
%     followup also enters), Cook et al. 2011, Dubovi/Levy/Dagan 2017.
%   M5.1 = 7 ¶¶ after pass. APPROVED 2026-07-28 — PASS COMPLETE.
%     superiority. FUTURE ITEM BANKED: expanded task-lever discussion w/
%     secondary pedagogical literature on proven interactive-learning-
%     environment efficacy — placement TBD (raise at M5.3). (P3.1)
%     FIND-LEVER DOCTRINE CORRECTED (AUTHOR): mandated student DOES
%     bring a governing desire — pass → degree → life = the
%     FOR-THE-SAKE-OF-WHICH (Heidegger tie rendered); routes vary
%     beneath it; YouTube-gravitation rational under route-parity;
%     design aim = ENGINEERED AGENCY — VLE better past parity, YouTube
%     stays a real option (quick refresh, phone), free choice keeps
%     landing on the world; Part II curation formula reused as internal
%     echo ("handed a world, never a word"; author's "hack their soul"
%     doctrine — soul-line reserved for M5.3's ethics beat). (P3.2)
%     second-pair-of-eyes EXPLAINED (multiple activity streams, assigned
%     vantage points, reconcile-table-with-room prompts) — moved to the
%     meeting portion as the find×meeting bridge. (P3.3) WHILE-LEVER
%     CORRECTED (AUTHOR): procedures must NOT be shortened — edits
%     misrepresent the operation and can delete the very corners where
%     unmet needs sit; levers = SEGMENTATION (weekly assigned portions +
%     follow-up) and/or INTERVAL INTERACTIVITY (prepared sequences that
%     interrupt the drag and redirect attention); "the while is not
%     shortened; it is articulated."
%     v4 RULINGS APPLIED 2026-07-28 (P2/P3 second round) — M5.1 now
%     6 ¶¶ (¶2 list+door · ¶3 task-obstacle · ¶4 task-on-the-chain ·
%     ¶5 find · ¶6 meeting+while [¶6 = prior ¶4, no rulings, carried]):
%     (P2.1) footnote quote REPLACED w/ unused W25-R035/Q10 (student-
%     designed in-world group task; char-exact vs CSV 2026-07-28);
%     (P2.2/2.3) "question the third term asked" REFRAMED as the
%     OBSTACLE the deployment never solved; "pedagogically prosaic" +
%     "world justified by compulsion is not justified at all" REMOVED
%     (author: Plato's cave is still a world) → replaced by (a) Aristotle
%     seven-causes compulsion-painful doctrine — paraphrase in prose +
%     verbatim footnote "So all acts of concentration, strong effort,
%     and strain are necessarily painful; they all involve compulsion
%     and force, unless we are accustomed to them, in which case it is
%     custom that makes them pleasant" VERIFIED char-exact vs
%     rhetorical_ontology Rhetoric_(2014) Clean Copy pdf p. 32, cite
%     I.11 1370a12–14 [marginal "15" lands at the following clause —
%     range pinned conservatively]; (b) G&G candle-lighting grounding
%     (USER-ORDERED cross-part reference): start-button-as-first-lesson,
%     grip/trigger, "exploration over instruction" = Part II v4's own
%     phrase reused internally, paraphrase only, no quote; (c)
%     choice-architecture doctrine: assignments completable outside
%     (YouTube + course page) OR everything in one place in-world;
%     illusion-of-agency clause rendered ("more arranged than open...
%     out of conscious recognition"); one-roof convenience.
%     (P2.4) task-on-the-chain ¶ added: chain links specified
%     (perception → phantasia's rendering → committal judgment w/
%     emotion concretizing → act; vocabulary per M3.3-v3 canon);
%     INCORPORATION-ACHIEVED correction (design cannot manufacture it —
%     achievement is the student's; design makes it achievable);
%     "in every way, shape, and form" close per user wording.
%     (P3.1) for-the-sake-of-which ITALICIZED; (P3.2) definition = "the
%     end from which every decision, conscious or otherwise, is
%     perceived, evaluated, understood, and enacted" (user wording);
%     (P3.3) "Routes are what vary beneath it" → concrete pursuit-ways
%     sentence; (P3.4) "same grade" → "same end"; (P3.5) YouTube-open
%     cases = no-strong-computer + quick-verify (user wording).
%     M3.6 "One voice, both poles" FIXED → "both sides" (backup
%     .backups/20260728T091044/; pole ban reaches it per user ruling).
%     v5 RULINGS APPLIED 2026-07-28 (third round): "Aristotle explains
%     why"; "the fumbling the lighting process takes"; ¶4 swap — USER'S
%     CANONICAL FORMULA "Rhetorical incorporation is a mode of
%     being-in-the-virtual-world: it is something curated, a
%     harmonization between the student and the virtual world within
%     which they dwell" (user dictated "user"; rendered "student" for
%     register consistency — FLAG for confirmation); ¶5 curating-the-
%     find sentence BROKEN UP (parity condition → past-parity elements,
%     one sentence each: prompts/target, captions/citable, deliverable-
%     linked finds, the waiting meeting, one-more-reason close). TWO NEW
%     STANDING RULES MEMORIZED: chain-node walkthrough = fundamental
%     analytical form ([[feedback-chain-node-analysis-fundamental]] —
%     user "Note to System"); bare "incorporation" = Calleja ONLY
%     ([[feedback-rhetorical-incorporation-naming]] — swept, 1 fix).)

What remains is design, and the diagnosis dictates its form. A design
lever, in this analysis, is an \textit{archē}: Aristotle's fourth
sense of the word names ``[t]hat from which \ldots\ a thing first
arises, and from which the movement or the change naturally first
proceeds'' (\textit{Metaphysics} V.1, 1013a7--10)---the productive
source from which a change first comes. To call something a design lever is to make a two-part claim: that one
of the deployment's documented failures---the quarter of the class
stopped at the installer, the forty-seven students who chose YouTube
over the world, the group meetings that never convened---has a definite
productive source in the design itself, and that design can reach that
source and change it. The five levers that follow are ranked by the
station of the actualization chain each one works on: the way into the
world (A\textsubscript{0}), the act the world houses (A\textsubscript{3}
to A\textsubscript{4}), the object the seeking is for, the convening of
the WE, and the shape of the watching's time. Ranked this way, the
levers recapitulate the diagnosis: each names a place where the
deployment fell silent, and the productive source design must work on
there.

Five levers carry that claim, and each can be given a plain name: the
\textit{door}, the \textit{task}, the \textit{find}, the
\textit{meeting}, and the \textit{while}. The \textit{door} is the way
into the world---the download, the installation, the sign-in, the
machine all of it must run on. A quarter of the class never got past
it; for those students the world disclosed nothing, because it never
opened. Everything design does inside a world waits on this lever,
since a world is only as good as its way in: a browser-deliverable
build, Mac support, and entry without setup re-admit the students the
\textit{door} turned away. The students proposed the same remedy
themselves: ``Provide a cloud version of the VR. Meaning, provide the VR
setup and ready to go on a cloud machine provided by UCI. This will help
for inclusivity in the VR'' (S25-R059/Q9)---entry without local
installation, asked for in the name of exactly the students the
\textit{door} excluded.

The \textit{task} is what there is to do once inside, and its absence
is the obstacle the deployment never solved. As deployed, every
assignment the course gave could be completed without entering the
world; the watching was mandated, but nothing had to be done where the
watching happened---and students asked for exactly what was
missing.\footnote{One proposed it in detail: ``Maybe an interactive
surgical platform where a group can work together to dissect/do
surgery to further understand what a surgery consists of \ldots\ it
has steps so a group can virtually perform the surgery''
(W25-R035/Q10).} The field's founding account of what
three-dimensional learning environments are for named this obstacle in
advance. What such environments distinctively offer, Dalgarno and Lee
argue, are tasks: ``well-designed 3-D VLEs can enable learning tasks
that are not possible or not as effective in 2-D environments''
(``Learning Affordances'' 25)---which is what the student's surgical
platform proposes, almost
to the letter. The same paper issues the warning this deployment went
on to demonstrate: comparisons that hold the learning design constant
across a 2-D and a 3-D environment ``would be likely to only
demonstrate the trivial fact that if the unique affordances of 3-D
VLEs are not harnessed within the learning design, there will be
minimal unique learning benefits'' (25). The deployment ran that
comparison without meaning to---the same design, watch the footage, on
both routes---and the cohort returned the predicted verdict in its own
words: ``the youtube videos were enough'' (W26-R033/Q9). A second
proposal closes a loop the analysis left open. The student whose praise
stood over an emptiness they could not name---the divergence register's
culprit-less voice---returned, in another answer, naming the missing
thing exactly: ``I think there should be reflections in the VR
environment. While being immersed in the procedure is fascinating, having
reflections afterwards would allow us to look for the main takeaways from
these procedures and allow us to think critically about clinical needs''
(W26-R016/Q9). What the second-form echo could not locate, the same
student's design request supplies: an act of reflection, housed where the
watching happens. One way to
remove the obstacle is requirement: make
the work possible only inside the world, and the rooms will fill. A
compulsory world is still a world. What compulsion changes is how the
world is entered and felt, and Aristotle explains why: what is done under
compulsion runs contrary to nature, and for that reason it is
painful---felt as force unless custom has long since softened
it.\footnote{``So all acts of concentration, strong effort, and strain
are necessarily painful; they all involve compulsion and force, unless
we are accustomed to them, in which case it is custom that makes them
pleasant'' (\textit{Rhetorica} I.11, 1370a12--14).} The dissertation's
earlier analysis of a tutorial world, \textit{Gnomes \& Goblins},
shows the alternative in miniature. That experience also needed a
beginning---no world starts without one---and it designed its start
button as a first lesson: to enter, the player must light a candle,
and the fumbling the lighting process takes introduces the grip, the trigger,
and the world's whole philosophy of exploration over instruction. The
requirement did not disappear; it was made to teach. The deployment's
version of the same design is choice. Let the assignments remain
completable outside---the YouTube platform and the course page stay
open---and build the environment so that everything can also be done
in one place, inside the world, where the footage, the finds, the
notes, and the group already live. A student who enters on those terms
enters by their own act. The force Aristotle warns of drops out of the
equation---or, where the choice is more arranged than open, out of
conscious recognition, which for the experience amounts to the same---
and the world gains a convenience no other route offers: the course's
whole work gathered under one roof.

What the \textit{task} should be built toward, the actualization chain
specifies link by link. The chain traces how a world becomes an act:
perception takes the world in; \textit{phantasia} renders what was
taken into a determinate image; judgment commits to what the image
presents, and emotion concretizes in that commitment; desire, so
formed, completes itself in an act. A task built on the chain gives
every link its material inside the world: footage worth the taking-in;
prompts that give the forming image something determinate to resolve
on; judgments the deliverable asks the student to commit to and
defend; and an act---the need noted, captured, submitted---housed in
the same rooms where the seeing happened, so that nothing in the
world's workings sends the student elsewhere to complete what the
world began. Design cannot manufacture rhetorical incorporation.
Rhetorical incorporation is a mode of
\textit{being-in-the-virtual-world}: it is something curated, a
harmonization between the student and the virtual world within which
they dwell; what design does is make it achievable, clearing the way
at every link.
Build the task so, and the world's existence stops needing an
argument: in every way, shape, and form that matters to the student's
own end, the world is perceived as better than any alternative.

The task lever's pedagogical warrant is already established at
meta-analytic scale. Across health-professions education, training
that puts the learner's hands on the procedure---simulated, not
real---has been tested against no training at all in 609 studies
enrolling 35,226 trainees, with large pooled effects on knowledge and
skills; the evidence is so uniform that the reviewers close by
doubting the comparison is worth running again: ``with only 4\% of the
outcomes failing to show a benefit, we question the need for further
studies comparing simulation with no intervention'' (Cook et al.\
987). That meta-analysis is medium-agnostic---mannequins and cadavers
count alongside screens---so it grounds the task's warrant without
ranking routes; the finer work belongs to the closest analogue this
deployment has in the literature. Dubovi, Levy, and Dagan built the
kind of world the student's proposal asks for: a non-immersive desktop
simulation of a clinical procedure---medication administration, in a
virtual hospital ward---where nursing students assess the patient,
prepare the dose, administer it, and monitor what follows. Students
who trained in that world scored higher on the study's knowledge
measures than a matched cohort taught the same material by lecture, and
the advantage was still undiminished when the students were retested
five months later (23). The study's sharpest
finding is the watching-versus-doing distinction made empirical: of
ten students whose sessions were recorded, nine answered the procedure
worksheet correctly, and nine nonetheless selected the wrong
medication on their first attempt inside the world---errors that fell
to near zero by the second scenario (23--24). Knowing the procedure on
paper and being able to perform it are different attainments, and only
a world with a task in it exposes the difference and then closes it.
The authors' design principle is the task lever's own conclusion,
reached by measurement: ``ensuring that participants gain sufficient
control of the environment is an important condition for learning,''
for ``presence is not a unique attribute of the environment, but it is
induced by the fidelity of representation and the high degree of
interaction and user control'' (25). That last sentence reaches past
the task to the door and the rooms: nothing about a desktop world
forecloses presence; fidelity---to a lesser extent---and interactive
control and agency---to a much greater extent---produce it.

The \textit{find} is the object of desire, and here the diagnosis must
be exact, because the student under mandate is not desireless. Every
student arrives carrying a governing desire that organizes the rest:
to pass the course, for the sake of the degree, for the sake of the
life the degree serves. In the vocabulary this dissertation has
carried from Heidegger, that is the student's
\textit{for-the-sake-of-which}---the end from which every decision,
conscious or otherwise, is perceived, evaluated, understood, and
enacted. What varies beneath that end is the way it is pursued: which
platform to open, where to watch, whether to watch alone. Watching on
the YouTube platform and watching in the virtual hospital are, to that
governing desire, two routes to the same end, and if the routes look
equally effective, the more convenient one wins every time. Curating the \textit{find} therefore means shaping the world so
that the student's own weighing lands inside it. The first condition
is parity: the build must at least equal the YouTube platform in video
quality and in everyday conveniences, because no curated object
survives a worse player. Past parity come the elements only the world
can add. Prompts name, for each procedure, what a needs-finder watches
for, so the looking has a target. Captions and transcripts mark the
moments that matter, so the target can be found and cited. Each find a student captures should feed directly into the group's
graded deliverable---an unmet need identified in the footage becomes an
entry the innovation notebook and the final report are built from---so
that the desire for the find is not a new desire the course must create
but the student's existing desire for the grade, extended one step down
into the world. The meeting waits inside, worth the
\textit{door}'s inconvenience. Each element gives the governing desire
one more reason to choose this room. The choice
must remain real. The YouTube platform stays open for the quick second
viewing, for the student without a computer strong enough to run the
environment, and for the one who wants to jump back in quickly to
verify something about an operation; the aim is not to close the other
route but to make the world the one a free choice keeps landing on.
Agency itself becomes the pedagogical instrument: the student moved by
being handed a world, never a word.

The \textit{meeting} is co-disclosedness convened. The proximity voice
chat affordance exists and every student already belongs to a procedure
group; what design must produce is the convening itself---whatever
brings two students into the same room at the same hour. It is the one
lever whose product no video can imitate, and the \textit{find} can be
shaped to invite it. A procedure fills the frame with more than one
stream of activity---the surgeon's hands, the instruments passing, the
movements of the team around the table---and no single watcher attends
to all of it at once. A prompt that assigns two vantage points, or a
find that asks what happened at the table to be reconciled with what
happened around it, makes a second pair of eyes worth more than a
second viewing: the desirable object doubled into an occasion to meet.
The \textit{while} is the temporal frame---the same while that boredom
stretches. The length complaints ran through every term, but the cut
they imply would be the wrong lever. A procedure edited for pace is no
longer the procedure, and what an edit removes may be exactly the
unremarkable corner of the operating room where an unmet need sits;
the footage must stay whole. The watching does not have to. Segment
the procedures into portions assigned week by week, and set interactive
sequences at intervals inside the
watching---prepared moments that interrupt the drag, catch the
attention as it wanders, and redirect it to what matters. The students
asked for this articulation in nearly those terms. One tied the length
directly to lost focus: ``The long videos of the procedures in the VR
could probably be improved. I didn't really understand what was going on,
and the videos felt too long. If they were split into multiple parts it
would probably be easier to focus on the procedure'' (S25-R017/Q6).
Another proposed the weekly portioning with its diagnosis attached: ``I
think a good think would be to assign certain VR on video to watch weekly
so that all members could stay on track. because the main issues was
people were just waiting too last minute to watch the videos''
(S25-R037/Q6)---segmentation not only against the drag of the while but
against the procrastination a term of total freedom invites. The while is
not shortened; it is articulated. Session length, sequence, and rest
belong to the levers, not the logistics.

% AUTHOR FACT (2026-07-31, batch 5): an instructor-run Discord channel
%   existed for the whole class in ALL surveyed terms (also in the S26
%   syllabus). Context for the meeting-lever remedies: the meeting's
%   nearest competitor channel had institutional standing of its own.
% --- M5.2 Collapsing the activation energy of the WE (★ expanded lever.
%     5-¶ architecture per approved plan: ¶1 equilibrium · ¶2 synchrony
%     constraint · ¶3 rung 0 + low ladder · ¶4 ladder top w/ ignition
%     logic + bia guard · ¶5 motto. Vocabulary ruling in force
%     (technical terms kept; chronos = design material, allotted/
%     reserved). ¶¶1–2 DRAFTED 2026-07-28 — in review.
%     Batch verifications (2026-07-28): W25-R060/Q10 (class google
%     sheet of group entry times) + S25-R079/Q6 (scheduled online mixer
%     early in course; NB row also attests a personality test used for
%     grouping — course fact, unused) char-exact vs open_text.csv ·
%     Kreijns Pitfall 2 "the tendency to restrict social interaction to
%     educational interventions aimed at cognitive processes" verified
%     vs corpus/pedagogy PDF, journal p. 336 (pdf p. 2; same page as
%     Pitfall 1, cite form per M1.6) — deploys in batch 2 ·
%     chronos = Part I §1.1 established term (co-implicated w/ kinēsis),
%     term-use only, no quote. Floors/equilibrium facts inherited from
%     M4.3/M3.5 — no re-derivation; "micro-synchrony" named in ¶2 (term
%     feeds M7.1's minimum-dose hand-forward).
%     P1/P2 RULINGS APPLIED 2026-07-28: (P1.1) "a decade of platform
%     design" REMOVED (misread as the author's own 5-year platform
%     work) → "the solo path's frictions have long since been engineered
%     away by the platforms that host it"; (P1.2) "one student confirmed
%     convening in the three terms the class has run" + FLOOR-HONESTY
%     FOOTNOTE (likely more than one convened; survey built without a
%     direct did-your-group-meet item — foresight admission consistent
%     w/ M3.5 ¶1's; count = floor, never a measure); (P2.1) ¶2 opener
%     rewritten concrete ("Every remedy for the meeting draws on one
%     material: shared time") — NEW STANDING RULE memorized:
%     [[feedback-no-vague-paragraph-frames]] (no vague paragraph
%     openers/closers; check first/last sentences in isolation).
%     ¶¶1–2 APPROVED 2026-07-28. ¶¶3–5 DRAFTED 2026-07-28 — in review.
%     Batch-2 verifications: W25-R071/Q9 ledger reuse (M3.5) ·
%     W25-R060/Q10 + S25-R079/Q6 char-exact (this session) · Kreijns
%     Pitfall 2 char-exact p. 336 (named authors in prose per M1.6
%     form) · De Back et al. 5375–76 + van der Meer et al. 9–11 = M1.6
%     locked-bank claims, parenthetical et-al forms · W26-R037/Q10
%     "main pro" fragment RE-VERIFIED char-exact this session — TRUE
%     WORDING IS "The main pro to the platform is you can watch with
%     others" (analysis docs' "the main pro is you can watch with
%     others" = paraphrase; draft quotes the true span w/ [t] bracket) ·
%     Aristotle custom-clause re-quoted from the M5.1 ¶3 footnote span
%     (I.11 1370a12–14, same verified sentence). Ignition/bia guard per
%     M3.3 ¶5 seven-causes inherit (archē-outside vocabulary); "graded
%     once if ignition needs it"; motto rephrased per vocabulary ruling
%     ("reserve the least chronos," never "spend the minimum").
%     v2 RULINGS APPLIED 2026-07-28 — ¶¶3–5 REWRITTEN as ¶¶3–4 (M5.2 now
%     4 ¶¶): LADDER/RUNG METAPHOR REMOVED (implied hierarchy; replaced
%     by plain ordering "by how much shared time each asks" +
%     first/second/third/fourth). (P3.1) AUTHOR FACT BANKED: the course
%     page ALREADY DESCRIBES the co-watching functionality → "syllabus
%     line" dropped; awareness remedy sharpened to description-tried-
%     and-failed → DEMONSTRATION. (P3.2) traces explained concretely
%     (trace shows the group uses the space; cannot bring two students
%     in at the same hour; no sign of yesterday's visit creates today's
%     meeting). (P3.3 + P4.1) AUTHOR DESIGN REPLACES system-proposed
%     default hour (imposing = moot + removes agency): SCHEDULING
%     ASSIGNMENT — each group settles its weekly time via a
%     group-scheduling tool (user's example: whenavailable.com — kept
%     generic in prose) + MINUTES FOLLOW-UP assignment (brief document:
%     what watched/found/settled; occurrence graded, timeslot never).
%     (P4.2) dyad spelled out ("a pair of groupmates…two students, one
%     procedure, one shared hour—or the full procedure group").
%     (P4.3) bia failure-mode EXPLAINED plainly (term-after-term
%     minutes-only meeting = origin outside agent, felt as force,
%     compliance not practice; opposite sequence spelled out w/
%     W26-R037 main-pro experience). (P5.1) ¶5 DELETED; goal folded
%     into ¶4 close: "an available affordance turned into a consciously
%     selected practice" (user wording).)

Of the five levers, the \textit{meeting} asks the most of design,
because everything needed for it already existed and it lapsed anyway.
Co-watching is a coordination game. Two students who would each get
more from watching together can each still watch alone, and the solo
path asks nothing of them: no schedules aligned, no invitation risked,
no hour kept. The meeting asks all three, and it asks them of all
students in the same group simultaneously. A coordination game of that shape has two stable
outcomes---everyone convenes, or no one does---and which outcome a
cohort lands on is decided by the default. Left to itself, the default
decides for solitude: the solo path's frictions have long since been
engineered away by the platforms that host it, while the meeting's
demands stand untouched. What the collected data showed is therefore
an equilibrium, not a verdict on the students: co-watching solved in
principle---the affordance built, the groups assigned---and unsolved
in practice, one student confirmed convening in the three terms the
class has run.\footnote{It is likely that more than one student
convened in the space with their group. The survey, however, was built
without the foresight of a question that asked directly whether the
assigned group had used the environment to meet, so the collected data
cannot say how many groups actually did. The count is a floor---the
number of responses that happen to confirm a convening---never a
measure of how many occurred.} The lever's name
is exact. What must collapse is the activation energy of the meeting:
the demands that stand between an assigned group and its first shared
hour.

Every remedy for the meeting draws on one material: shared time. The
course was built, deliberately and for good reason, to need none of
it. Its students enroll across nine campuses, quarters set against
semesters, with no timetable in common; a synchronous version of the
course was never in the design space, and an in-person one still
less. Whole-cohort
synchrony is therefore closed to design. What remains open is small:
the overlap of two schedules, a procedure group's single shared
hour---micro-synchrony, the dyad-sized coincidence the systemwide
constraint still permits. Time, in this frame, is a design
material. The deployment allotted all of it to access---any student,
any hour, any place---and that allotment was right; the error was only
that nothing was held back. The meeting needs a small reserve of
shared time, and the deployment reserved none. The design problem is
exactly that narrow: not to restore the classroom hour the course
simply could not implement, but to reserve, inside a schedule built on
total freedom, the one small overlap the WE requires.
% RESOLVED 2026-07-31 (user, close-out): semester-UC question verified by
% the author — the batch-5 facts stand (open to Berkeley/Merced; no
% semester-campus enrollees thus far; rendered in M0 ¶6).

The remedies can be ordered by how much shared time each asks of the
group, and the first asks none: demonstrate the capability. Twelve of
the thirteen students who asked for a way to watch together asked in
wording that implied the platform had none---``if there was a way to
interact with other users who are online at the same time''
(W25-R071/Q9)---and this despite the course page, which states plainly
that groups can meet and talk inside the environment. The capability
was described in writing and asked for as if absent; what description
did not accomplish, demonstration can---an onboarding walkthrough, a
live demonstration in the first week's module, the capability seen in
use rather than read about. The second remedy also asks no shared
time: let the world carry traces of the group---who entered today,
which rooms a groupmate visited. A trace shows a student that their
group uses the space; what it cannot do is bring two students into the
space at the same hour. Watching together by proximity voice chat happens
only when both students are present at once, and no sign of
yesterday's visit creates today's meeting; traces point toward the
meeting without producing it. The third remedy is an assignment, and
it produces the meeting's hour while leaving the choice of hour
entirely with the group. Early in the term, each procedure group uses
a group-scheduling tool to settle the weekly time its members can
share; the time is the group's own finding, not the course's
imposition. The students suggested this remedy directly themselves:
one proposed ``setting up a class google sheet where times are set by
each group to enter the VR program/watch the videos'' (W25-R060/Q10);
another, a ``scheduled online mixer or networking session early in the
course'' to make finding groupmates easier (S25-R079/Q6). The second
proposal reaches past scheduling into affiliation, and the
collaborative-learning literature says it must: Kreijns, Kirschner,
and Jochems name as their second pitfall ``the tendency to restrict
social interaction to educational interventions aimed at cognitive
processes'' (336)---a meeting made only of assignments starves the
social ground the meeting stands on. The students also supplied a reason
the world's own meeting-form can carry that social ground better than a
video call: ``I think a virtual group will be nice, students could make
their own charactors, students who don't know each other may shy to speak
in reality, but if there is a virtual image for them, they might have
less stress talking to others'' (W26-R001/Q10). Two more voices
corroborate the mechanism from the camera's side, asking for ``group
communications that do not wish to do video camera meetings''
(S25-R002/Q10) and noting that ``some of us don’t wanna have their
cameras on'' (S25-R074/Q10). The avatar lowers the exposure the
invitation risks; for strangers assigned to one another across nine
campuses, that reduction works on exactly the activation energy this
lever must collapse.

The fourth remedy closes the design with accountability, still without
imposing a timeslot. Each group's weekly meeting carries a brief
follow-up assignment: minutes, submitted after the session---what the
group watched, what it found, what it settled. The occurrence is
graded; the timeslot never is; the hour remains the group's own
settlement from the scheduling assignment. Small meetings are the
right form. In immersed collaborative environments, learning gains run
inverse to group size---smaller groups, larger gains (De Back et al.\
5375--76)---and reviews of monitor-based collaborative worlds judge
voice communication essential (van der Meer et al.\ 9--11);
groupmates chatting through proximity voice chat---two students, one
procedure, one shared hour---or the full procedure group at once, is
exactly the form the evidence favors. A requirement stands at the start of this
design, and Aristotle's qualification, quoted above, explains why it
can stand there without poisoning what follows: the forced is painful
``unless we are accustomed'' to it, and then ``it is custom that
makes'' it pleasant. The assignment's work is to establish and naturalize the custom---to
bring each group through its first meetings, while the demands are
still unfamiliar and the worth still unexperienced. What it curates
is, in Aristotle's terms, a habit: among the seven causes of action,
the one that stands nearest to nature, for ``as soon as a thing has
become habitual, it is virtually natural'' (\textit{Rhetorica} I.11,
1370a6--8). If that were all
the design ever did---if, term after term, students met only because
minutes were due---the meeting would have remained what the
seven-causes analysis names compulsion: an act whose origin stays
outside the agent, done under force and felt as force, yielding
compliance rather than practice. The design intends the opposite
sequence. The requirement produces the first meeting; the first
meeting lets the students experience what the survey's one confirmed
convener reported---``[t]he main pro to the platform is you can watch
with others'' (W26-R037/Q10)---and from then on the group meets
because it wants to, the cause now inside the students rather than in
the requirement. The goal, throughout, is exactly that movement: an
available affordance turned into a consciously selected practice. That
movement is also a discovery the actualization chain makes diagnosable.
Requirement and practice can look identical from the outside---the same
group, the same room, the same hour---and the chain names where they
differ: in the \textit{arch\=e} of the act, its origin lying in the
requirement or in the students' own desire. A course that watches only
attendance cannot tell compliance from practice; a course that asks
where the act's origin now lies can see, term by term, whether the
affordance is still being forced or has begun to be chosen.

% --- M5.3 The motto and the mirror (1 ¶ per approved plan. DRAFTED
%     2026-07-28 — in review. Verifications: Mark capture-economy claim
%     = paraphrase, cite (Mark 155–65) per crosswalk §3 catalog-backed
%     range — VERIFY Mark is in Works Cited at assembly (M0/M1 may
%     carry); Wendt "facilitate profitable mistakes" char-exact vs
%     corpus/new_media Design for Dasein PDF p. 162 (pdf 164; OCR noise
%     in layer, wording confirmed) — "profitable" appears ONLY inside
%     Wendt's quoted phrase, not our own metaphor; Wendt principle 10
%     ("advocate for preferable worlds; refuse dark patterns") = entry
%     manual's distillation, deployed as paraphrase. Soul-line = Part II
%     v4's "reaching into and reshaping the soul" reused as internal
%     formula (the author's "hack their soul" doctrine landing at the
%     ethics beat as reserved). "Settled ways" not "habits" (hexis/habit
%     ban). M5.2's for-the-sake-of-which + capture/restart mirror from
%     cross-case §3.3.
%     v2 RULINGS APPLIED 2026-07-28: (1) attention-economy passage
%     EXPANDED w/ Mark catalog verbatims — "who you are, how you feel,
%     what you do when and where" + "to capture your attention" (155,
%     item #28); "occur spontaneously and automatically without
%     cognitive processing" + "when you are low in attentional resources
%     or in a bored or a rote state" + "you can't help but respond to
%     them" (159–60, items #34–35; all MARK-VERBATIM register).
%     (2) NO-STALL DOCTRINE (user): the chain NEVER stalls — "stalled
%     chain" lines DELETED; capture rendered as full chain walkthrough
%     (no determinate object → pre-charged object supplied → perception
%     seized → act completes past the cognitive orientations, "without
%     cognitive processing" = Mark's own warrant); "Nothing in the chain
%     is broken; the chain runs. It is run deliberately past the very
%     links where a person's own ends would get a vote." Doctrine added
%     to [[feedback-chain-node-analysis-fundamental]] + "stall" added to
%     sweeps; M5.1 ¶1 "a stall" → "a flight" (approved-text touch).
%     (3) DARK-PATTERNS USE CORRECTED (user suspicion CONFIRMED vs PDF:
%     Wendt lists dark patterns as one nefarious effect — "persuasive
%     design attempting to directly influence behavior--dark patterns
%     attempting to facilitate profitable mistakes" — NOT a name for
%     attention capture): capture-equation dropped; Wendt now supplies
%     the designer's norm, "advocating for preferable worlds" VERIFIED
%     char-exact p. 162 (pdf 164). (4) M5.2 ¶4 HABIT INSERTION:
%     Aristotle's habit curated — "as soon as a thing has become
%     habitual, it is virtually natural" VERIFIED char-exact pdf p. 31
%     (Rhetoric I.11 1370a6–8, M3.3's cite; habit = seven-causes ethos,
%     kept distinct from hexis per standing ban). (5) soul-line = "To
%     design this way is to hack the soul." (user wording).
%     v3 2026-07-28: "video alone cannot" (user). ¶ APPROVED —
%     **M5.3 COMPLETE. M5 MODULE COMPLETE (M5.1 6¶¶ · M5.2 4¶¶ ·
%     M5.3 1¶ — all approved). M5.2 habit insertion approved.**)

The five levers share one motto: design against the question ``why not
just YouTube?'' Every element of the world must do something video
alone cannot---give a task, hold a place, convene a meeting---or the
element loses, rationally, to the platform that plays the same footage
with less friction. A craft this effective deserves candor about what it
is, because the attention economy practices the same craft in mirror
image, and the actualization chain shows exactly where the mirror
sits. Gloria Mark's account of attention capture is precise about the
mechanism. The algorithms learn ``who you are, how you feel, what you
do when and where'' and use it ``to capture your attention'' (Mark
155). Their instruments are targeted solicitations built to elicit
what Mark calls lower-order emotions---reactions that ``occur
spontaneously and automatically without cognitive processing''---and
they are timed for the moments ``when you are low in attentional
resources or in a bored or a rote state,'' when ``you can't help but
respond to them'' (159--60). Read through the chain, the design is
exact. A bored viewer still desires---the governing end never
sleeps---but the moment offers no determinate object to want. The
application supplies one, pre-charged: perception is seized by the
notification; the image arrives already pleasurable; and the act
completes before the cognitive orientations have done their
work---no deliberative weighing of the object, no committal judgment
about its worth. Nothing in the chain is broken by such design; the
chain runs. The pedagogical craft uses the same knowledge in the
opposite direction: it curates a worthy object of desire, informs the
judgment that weighs it, and gives committed desire an act to complete
in the world---every link furnished, none bypassed, the completion
serving the student's own \textit{for-the-sake-of-which} rather than a
platform's. Wendt states the designer's side of the same boundary:
what designers make acts on the world long after they are done with
it, the nefarious end of the spectrum runs from persuasive design to
dark patterns, and the calling that stands against it is
``advocating for preferable worlds'' (\textit{Design for Dasein} 162).
To design this way is to hack the soul. That is no reason to refuse
the craft; it is the reason the craft bears the responsibility it
does, and must answer for the dispositions it lays down---the settled
ways of looking, judging, and meeting that outlast the course. A world
that teaches by being dwelt in shapes its dwellers, and the designer's
obligation is to shape them toward their own good---openly aimed, and
defensible to the very students it moves.
